\newcommand{\titulus}{\nomenFesti{S. Hieronymi, Presbyteri \& Ecclesiæ Doctoris.}
\dies{Die 30. Septembris.}}
\newcommand{\oratio}{\pars{Oratio.}

\noindent Deus, qui beáto Hierónymo, presbýtero, suávem et vivum Scriptúræ sacræ afféctum tribuísti, da, ut pópulus tuus verbo tuo ubérius alátur et in eo fontem vitæ invéniat.

\pars{Pro pace in universo mundo.} \scriptura{Sir. 50, 25; 2 Esdr. 4, 20; \textbf{H416}}

\vspace{-4mm}

\antiphona{II D}{temporalia/ant-dapacemdomine.gtex}

\vfill

\noindent Deus, a quo sancta desidéria, recta consília et iusta sunt ópera: da servis tuis illam, quam mundus dare non potest, pacem; ut et corda nostra mandátis tuis dédita, et hóstium subláta formídine, témpora sint tua protectióne tranquílla.

\noindent Per Dóminum nostrum Iesum Christum, Fílium tuum, qui tecum vivit et regnat in unitáte Spíritus Sancti, Deus, per ómnia sǽcula sæculórum.

\noindent \Rbardot{} Amen.}
\newcommand{\invitatorium}{\pars{Invitatorium.}

\vspace{-4mm}

\antiphona{IV}{temporalia/inv-fontemsapientiae.gtex}}
\newcommand{\hymnusmatutinum}{\pars{Hymnus}

\cuminitiali{IV}{temporalia/hym-AEterneSol.gtex}}
\newcommand{\matversus}{\noindent \Vbardot{} Omnes mirabántur in verbis grátiæ.

\noindent \Rbardot{} Quæ procedébant de ore ipsíus.}
\newcommand{\lectioi}{\pars{Lectio I.} \scriptura{Phil. 2, 12-30}

\noindent De Epístola beáti Pauli apóstoli ad Philippénses.

\noindent Fratres Caríssimi mei, sicut semper obœdístis, non ut in præséntia mei tantum sed multo magis nunc in abséntia mea, cum metu et tremóre vestram salútem operámini; Deus est enim, qui operátur in vobis et velle et perfícere pro suo beneplácito. Omnia fácite sine murmuratiónibus et hæsitatiónibus, ut efficiámini sine queréla et símplices, fílii Dei sine reprehensióne in médio generatiónis pravæ et pervérsæ, inter quos lucétis sicut luminária in mundo, verbum vitæ fírmiter tenéntes ad glóriam meam in die Christi, quia non in vácuum cucúrri neque in vácuum laborávi. Sed et si delíbor supra sacrifícium et obséquium fídei vestræ, gáudeo et congáudeo ómnibus vobis; idípsum autem et vos gaudéte et congaudéte mihi.

\noindent Spero autem in Dómino Iesu Timótheum cito me míttere ad vos, ut et ego bono ánimo sim, cógnitis, quæ circa vos sunt. Néminem enim hábeo tam unánimem, qui sincére pro vobis sollícitus sit; omnes enim sua quærunt, non quæ sunt Iesu Christi. Probatiónem autem eius cognóscitis, quóniam sicut patri fílius mecum servívit in evangélium. Hunc ígitur spero me míttere, mox ut vídero, quæ circa me sunt; confído autem in Dómino, quóniam et ipse cito véniam.

\noindent Necessárium autem existimávi Epaphrodítum fratrem et cooperatórem et commilitónem meum, vestrum autem apóstolum et minístrum necessitátis meæ, míttere ad vos, quóniam omnes vos desiderábat et mæstus erat, proptérea quod audierátis illum infirmátum. Nam et infirmátus est usque ad mortem, sed Deus misértus est eius; non solum autem eius, verum et mei, ne tristítiam super tristítiam habérem. Festinántius ergo misi illum, ut, viso eo, íterum gaudeátis, et ego sine tristítia sim. Excípite ítaque illum in Dómino cum omni gáudio et eiúsmodi cum honóre habetóte, quóniam propter opus Christi usque ad mortem accéssit in intéritum tradens ánimam suam, ut suppléret id, quod vobis déerat ministérii erga me.}
\newcommand{\responsoriumi}{\pars{Responsorium 1.} \scriptura{\Rbardot{} Ps. 60, 2-3 \Vbardot{} ibid., 3-4; \textbf{H87}}

\vspace{-5mm}

\responsorium{VII}{temporalia/resp-exaudideusdeprecationem-CROCHU.gtex}{}}
\newcommand{\lectioii}{\pars{Lectio II.} \scriptura{Nn. 1.2: CCL 73, 1-3}

\noindent Ex Prólogo commentariórum sancti Hierónymi presbýteri in Isaíam prophétam.

\noindent Reddo quod débeo, obœ́diens Christi præcéptis, qui ait: \textit{Scrutámini Scriptúras;} et: \textit{Quǽrite et inveniétis,} ne illud áudiam cum Iudǽis: \textit{Errátis, nesciéntes Scriptúras neque virtútem Dei.} Si enim iuxta apóstolum Paulum, Christus Dei virtus est Deíque sapiéntia, et qui nescit Scriptúras nescit Dei virtútem eiúsque sapiéntiam: ignorátio Scripturárum ignorátio Christi est.

\noindent Unde imitábor patremfamílias, qui de thesáuro suo profert nova et vétera; et sponsam dicéntem in Cántico canticórum: \textit{Nova et vétera, fratruélis meus, servávi tibi;} sicque expónam Isaíam, ut illum non solum prophétam, sed evangelístam et apóstolum dóceam. Ipse enim de se et de céteris evangelístis ait: \textit{Quam speciósi pedes evangelizántium bona, evangelizántium pacem.} Et ad ipsum quasi ad apóstolum lóquitur Deus: \textit{Quem mittam, et quis ibit ad pópulum istum?} Et ille respóndit: \textit{Ecce ego, mitte me.}}
\newcommand{\responsoriumii}{\pars{Responsorium 2.} \scriptura{\textbf{H373}}

\vspace{-5mm}

\responsorium{III}{temporalia/resp-istecognovit-CROCHU.gtex}{}}
\newcommand{\lectioiii}{\pars{Lectio III.}

\noindent Nullúsque putet me volúminis istíus arguméntum brevi cúpere sermóne comprehéndere, cum univérsa Dómini sacraménta præsens Scriptúra contíneat; et tam natus de Vírgine Emmánuel, quam illústrium patrátor óperum atque signórum, mórtuus ac sepúltus, et resúrgens ab ínferis, et Salvátor universárum géntium prædicétur. Quid loquar de phýsica, éthica et lógica? Quidquid sanctárum est Scripturárum, quidquid potest humána lingua proférre et mortálium sensus accípere, isto volúmine continétur. De cuius mystériis testátur ipse qui scripsit: \textit{Et erit vobis vísio ómnium, sicut verba libri signáti, quem cum déderint sciénti lítteras, dicent: Lege istum. Et respondébit: Non possum, signátus est enim. Et dábitur liber nesciénti lítteras dicetúrque ei: Lege. Et respondébit: Néscio lítteras.}

\noindent Quod si cui vidétur infírmum, illud eiúsdem Apóstoli áudiat: \textit{Prophétæ duo aut tres loquántur, et álii diiúdicent; si autem álii fúerit revelátum sedénti, prior táceat.} Qua possunt ratióne reticére, cum in dicióne sit Spíritus qui lóquitur per prophétas, vel tacére vel dícere? Si ergo intellegébant quæ dicébant, cuncta sapiéntiæ rationísque sunt plena. Nec aer voce pulsátus ad aures eórum perveniébat; sed Deus loquebátur in ánimo prophetárum, iuxta illud quod álius Prophéta dicit: \textit{Angelus qui loquebátur in me}, et: \textit{ Clamántes in córdibus nostris, Abba, Pater}, et: \textit{ Audiam quid loquátur in me Dóminus Deus.}}
\newcommand{\responsoriumiii}{\pars{Responsorium 3.} \scriptura{\Rbardot{} Sir. 15, 3 \Vbardot{} Ps. 131, 18; \textbf{H62}}

\vspace{-5mm}

\responsorium{VII}{temporalia/resp-cibavitillum-CROCHU-cumdox.gtex}{}}
\newcommand{\hymnuslaudes}{\pars{Hymnus}

\cuminitiali{VIII}{temporalia/hym-FestivaCanimus.gtex}}
\newcommand{\lectiobrevis}{\pars{Lectio Brevis.} \scriptura{Sap. 7, 13-14}

\noindent Sapiéntiam sine fictióne dídici et sine invídia commúnico; divítias illíus non abscóndo. Infinítus enim thesáurus est homínibus; quem qui acquisiérunt, ad amicítiam in Deum se paravérunt propter disciplínæ dona commendáti.}
\newcommand{\responsoriumbreve}{\pars{Responsorium breve.} \scriptura{Eccli. 44, 15}

\cuminitiali{VI}{temporalia/resp-sapientiamsanctorum.gtex}}
\newcommand{\preces}{\noindent Christo, bono pastóri,~\gredagger{} qui pro suis óvibus ánimam pósuit,~\grestar{} laudes grati exsolvámus et supplicémus, dicéntes:

\Rbardot{} Pasce pópulum tuum, Dómine.

\noindent Christe, qui in sanctis pastóribus misericórdiam et dilectiónem tuam dignátus es osténdere,~\grestar{} numquam désinas per eos nobíscum misericórditer ágere.

\Rbardot{} Pasce pópulum tuum, Dómine.

\noindent Qui múnere pastóris animárum fungi per tuos vicários pergis,~\grestar{} ne destíteris nos ipse per rectóres nostros dirígere.

\Rbardot{} Pasce pópulum tuum, Dómine.

\noindent Qui in sanctis tuis, populórum dúcibus, córporum animarúmque médicus exstitísti,~\grestar{} numquam cesses ministérium in nos vitæ et sanctitátis perágere.

\Rbardot{} Pasce pópulum tuum, Dómine.

\noindent Qui, prudéntia et caritáte sanctórum, tuum gregem erudísti,~\grestar{} nos in sanctitáte iúgiter per pastóres nostros ædífica.

\Rbardot{} Pasce pópulum tuum, Dómine.}
\newcommand{\benedictus}{\pars{Canticum Zachariæ.} \scriptura{Sir. 15, 3; \textbf{H61}}

\vspace{-4mm}

\antiphona{VII c}{temporalia/ant-cibaviteumdominus.gtex}

\vspace{-3mm}

\scriptura{Lc. 1, 68-79}

\vspace{-2mm}

\cantusSineNeumas
\initiumpsalmi{temporalia/benedictus-initium-vii-c-auto.gtex}

\vspace{-1.5mm}

\input{temporalia/benedictus-vii-c.tex} \Abardot{}}
\newcommand{\benedicamuslaudes}{\cuminitiali{}{temporalia/benedicamus-memoria-laudes.gtex}}
\newcommand{\hebdomada}{infra Hebdom. XXVI per Annum.}
\newcommand{\matub}{Matutinum Hebdomadae B}
\newcommand{\matubd}{Matutinum Hebdomadae B vel D}
\newcommand{\laudb}{Laudes Hebdomadae B}
\newcommand{\laudbd}{Laudes Hebdomadae B vel D}

% LuaLaTeX

\documentclass[a4paper, twoside, 12pt]{article}
\usepackage[latin]{babel}
%\usepackage[landscape, left=3cm, right=1.5cm, top=2cm, bottom=1cm]{geometry} % okraje stranky
%\usepackage[landscape, a4paper, mag=1166, truedimen, left=2cm, right=1.5cm, top=1.6cm, bottom=0.95cm]{geometry} % okraje stranky
\usepackage[landscape, a4paper, mag=1400, truedimen, left=0.5cm, right=0.5cm, top=0.5cm, bottom=0.5cm]{geometry} % okraje stranky

\usepackage{fontspec}
\setmainfont[FeatureFile={junicode.fea}, Ligatures={Common, TeX}, RawFeature=+fixi]{Junicode}
%\setmainfont{Junicode}

% shortcut for Junicode without ligatures (for the Czech texts)
\newfontfamily\nlfont[FeatureFile={junicode.fea}, Ligatures={Common, TeX}, RawFeature=+fixi]{Junicode}

\usepackage{multicol}
\usepackage{color}
\usepackage{lettrine}
\usepackage{fancyhdr}

% usual packages loading:
\usepackage{luatextra}
\usepackage{graphicx} % support the \includegraphics command and options
\usepackage{gregoriotex} % for gregorio score inclusion
\usepackage{gregoriosyms}
\usepackage{wrapfig} % figures wrapped by the text
\usepackage{parcolumns}
\usepackage[contents={},opacity=1,scale=1,color=black]{background}
\usepackage{tikzpagenodes}
\usepackage{calc}
\usepackage{longtable}
\usetikzlibrary{calc}

\setlength{\headheight}{14.5pt}

% Commands used to produce a typical "Conventus" booklet

\newenvironment{titulusOfficii}{\begin{center}}{\end{center}}
\newcommand{\dies}[1]{#1

}
\newcommand{\nomenFesti}[1]{\textbf{\Large #1}

}
\newcommand{\celebratio}[1]{#1

}

%%%% Prejmenovat na latinske
\newcounter{Thumb}
\setcounter{Thumb}{0}
\newcommand{\thumbTitulus}{}

\newif\ifSideThumbsEnabled

\newlength\ThumbSize
\setlength\ThumbSize{2cm}

\SideThumbsEnabledfalse
\AddEverypageHook{
\ifSideThumbsEnabled
\ifodd\value{page}
\backgroundsetup{
angle=90,
position={current page.east|-current page text area.north east},
vshift=12pt,
hshift=-\value{Thumb}*\ThumbSize+\ThumbSize+\ThumbSize/4,
contents={\tikz\node[fill=gray!30,anchor=west,text width=\ThumbSize,
 align=center,text height=14pt,text depth=30pt]
{\small\thumbTitulus};
}
}
\else
\backgroundsetup{
angle=90,
position={current page.west|-current page text area.north west},
vshift=-16pt,
hshift=-\value{Thumb}*\ThumbSize+\ThumbSize+\ThumbSize/4,
contents={\tikz\node[fill=gray!30,anchor=west,text width=\ThumbSize,
 align=center,text height=28pt,text depth=6pt]
{\small\thumbTitulus};
}
}
\fi
\BgMaterial
\else\relax\fi}
\newcommand\sideThumbs[1]{\SideThumbsEnabledtrue%
\renewcommand{\thumbTitulus}{#1}%
\stepcounter{Thumb}
}
\newcommand\RemoveSideThumbs{\SideThumbsEnabledfalse}

\newcommand{\hora}[1]{%
\vspace{0.5cm}{\large \textbf{#1}}

\fancyhead[LE]{\thepage\ / #1}
\fancyhead[RO]{#1 / \thepage}
\addcontentsline{toc}{subsection}{#1}
}

% In Junicode, the default R/ with .35em has the bar too close to the
% letter.
\grelatexsimpledefbarredsymbol{A}{.30em}{.25em}{.30em}{.27em}%
\grelatexsimpledefbarredsymbol{R}{.41em}{.45em}{.45em}{.48em}%
\grelatexsimpledefbarredsymbol{V}{.10em}{.10em}{.10em}{.13em}%

% Emit R/, V/, A/ and + in red
\let\oldrbar\Rbar
\let\oldvbar\Vbar
\let\oldabar\Abar
\let\oldgrecross\grecross
\let\oldgrealtcross\grealtcross
\let\oldGreStar\GreStar
\let\oldGreDagger\GreDagger
\let\olddag\dag
\let\oldddag\ddag
\renewcommand{\Rbar}{{\color{red}\oldrbar}}
\renewcommand{\Vbar}{{\color{red}\oldvbar}}
\renewcommand{\Abar}{{\color{red}\oldabar}}
\renewcommand{\grecross}{{\color{red}\oldgrecross}}
\renewcommand{\grealtcross}{{\color{red}\oldgrealtcross}}
\newcommand{\grestar}{{\color{red}\oldGreStar}}
\renewcommand{\GreStar}{{\color{red}\oldGreStar}}
\newcommand{\gredagger}{{\color{red}\oldGreDagger}}
\renewcommand{\GreDagger}{{\color{red}\oldGreDagger}}
\renewcommand{\dag}{{\color{red}\olddag}}
\renewcommand{\ddag}{{\color{red}\oldddag}}
\newcommand{\Rbardot}{{\color{red}\oldrbar.~}}
\newcommand{\Vbardot}{{\color{red}\oldvbar.~}}
\newcommand{\Abardot}{{\color{red}\oldabar.~}}
\newcommand{\Psdot}{{\color{red}Ps.~}}

\let\quasiHora\hora

% larger unit than a hora
\newcommand{\divisio}[1]{%
\begin{center}
{\Large \textsc{#1}}
\end{center}
\fancyhead[CO,CE]{#1}
\addcontentsline{toc}{section}{#1}
}

% a part of a hora, larger than pars
\newcommand{\subhora}[1]{%
\begin{center}
{\large \textit{#1}}
\end{center}
%\fancyhead[CO,CE]{#1}
\addcontentsline{toc}{subsubsection}{#1}
}

% rubricated inline text
\newcommand{\rubricatum}[1]{\textit{#1}}

% standalone rubric
\newcommand{\rubrica}[1]{\vspace{3mm}\rubricatum{#1}}

\newcommand{\notitia}[1]{\textcolor{red}{#1}}

\newcommand{\scriptura}[1]{\hfill \small\textit{#1}}

\newcommand{\translatioCantus}[1]{\vspace{1mm}%
{\noindent\footnotesize \nlfont{#1}}}

\newenvironment{translatioMulticol}[1]{%
  \begin{multicols}{#1}
    \footnotesize
    \setlength{\parindent}{0pt}
    \setlength{\columnseprule}{0pt}
}{\end{multicols}}

% pruznejsi varianta nasledujiciho - umoznuje nastavit sirku sloupce
% s prekladem
\newcommand{\psalmusEtTranslatioB}[3]{%
  \vspace{0.5cm}
  \begin{parcolumns}[colwidths={2=#3}, nofirstindent=true]{2}
    \colchunk{%
      \input{#1}
    }

    \colchunk{%
      \vspace{-0.5cm}
      {\footnotesize \nlfont
        \input{#2}
      }
    }
  \end{parcolumns}
}

\newcommand{\psalmusEtTranslatio}[2]{%
  \psalmusEtTranslatioB{#1}{#2}{8.5cm}
}

\newcommand{\psalmusTr}[1]{{\footnotesize\nlfont\hspace*{-4pt} #1}}
\newcommand{\psalmusEtTranslatioT}[2]{{%
  \setlength{\LTleft}{-0.25cm}
  \setlength{\LTright}{\fill}
  \setlength{\tabcolsep}{0.35cm}
  \begin{psalmus}
      \noindent
      \begin{longtable}{p{\textwidth-#2-1cm}p{#2}}
        \input{#1}
      \end{longtable}
  \end{psalmus}
}}
\newcommand{\psalmusEtTranslatioTS}[2]{{%
  \setlength{\LTleft}{-0.25cm}
  \setlength{\LTright}{\fill}
  \setlength{\tabcolsep}{0.25cm}
  \begin{psalmus}
      \noindent
      \begin{longtable}{p{\textwidth-#2-0.5cm}p{#2}}
        \input{#1}
      \end{longtable}
  \end{psalmus}
}}

\newcommand{\canticumMagnificatEtTranslatio}[1]{%
  \psalmusEtTranslatioB{#1}{temporalia/extra-adventum-vespers/magnificat-boh.tex}{12cm}
}
\newcommand{\canticumBenedictusEtTranslatio}[1]{%
  \psalmusEtTranslatioB{#1}{temporalia/extra-adventum-laudes/benedictus-boh.tex}{10.5cm}
}

\newcommand{\textusEtTranslatio}[3]{%
  \begin{parcolumns}[colwidths={2=#3}, nofirstindent=true]{2}
    \colchunk{%
      #1
    }

    \colchunk{%
      \vspace{-0.5cm}
      {\footnotesize \nlfont
        #2
      }
    }
  \end{parcolumns}
}

% volne misto nad antifonami, kam si zpevaci dokresli neumy
\newcommand{\hicSuntNeumae}{\vspace{0.5cm}}

% prepinani mista mezi notovymi osnovami: pro neumovane a neneumovane zpevy
\newcommand{\cantusCumNeumis}{%
  \grechangestaffsize{14}
%  \global\advance\grespaceabovelines by 5mm%
%  \global\advance\grespaceabovelines by 2mm%
}
\newcommand{\cantusSineNeumas}{%
  \grechangestaffsize{14}
%  \global\advance\grespaceabovelines by -5mm%
  \grechangedim{spaceabovelines}{0cm}{scalable}%
}

\grechangestyle{abovelinestext}{\color{red}\it}

% znaky k umisteni nad inicialu zpevu
\newcommand{\superInitialam}[1]{\greannotation{\footnotesize {\textbf{#1}}}}

% pars officii, i.e. "oratio", ...
\newcommand{\pars}[1]{\textbf{#1}}

\newenvironment{psalmus}{%
  \setlength{\parindent}{0pt}
  \setlength{\parskip}{5pt}
}{%
  \setlength{\parindent}{10pt}
  \setlength{\parskip}{10pt}
}

%%%% Prejmenovat na latinske:
\newcommand{\nadpisZalmu}[1]{%
  \hspace{2cm}\textbf{#1}\vspace{2mm}%
  \nopagebreak%

}

% mode, score, translation
\newcommand{\antiphona}[3]{%
  \hicSuntNeumae%
  \superInitialam{#1}%
  \gresetinitiallines{1}%
  \gregorioscore{#2}%

#3
}

% a semantic alias of the previous
\newcommand{\responsorium}[3]{\antiphona{#1}{#2}{#3}}

\newcommand{\initiumpsalmi}[1]{%
  \gresetinitiallines{0}%
  \gregorioscore{#1}%
}

\newcommand{\cuminitiali}[2]{%
  \superInitialam{#1}%
  \gresetinitiallines{1}%
  \gregorioscore{#2}%
}

\newcommand{\sineinitiali}[1]{%
  \gresetinitiallines{0}%
  \gregorioscore{#1}%
}
 % Often used macros

\newcommand{\annusEditionis}{2023}

%%%% Vicekrat opakovane kousky

\newcommand{\anteOrationem}{
  \rubrica{Ante Orationem, cantatur a Superiore:}

  \pars{Supplicatio Litaniæ.}

  \cuminitiali{}{temporalia/supplicatiolitaniae.gtex}

  \pars{Oratio Dominica.}

  \cuminitiali{}{temporalia/oratiodominica.gtex}

  \rubrica{Deinde dicitur ab Hebdomadario:}

  \cuminitiali{}{temporalia/dominusvobiscum-solemnis.gtex}

  \rubrica{In choro monialium loco Dominus vobiscum dicitur:}

  \sineinitiali{temporalia/domineexaudi.gtex}
}

\setlength{\columnsep}{30pt} % prostor mezi sloupci

%%%%%%%%%%%%%%%%%%%%%%%%%%%%%%%%%%%%%%%%%%%%%%%%%%%%%%%%%%%%%%%%%%%%%%%%%%%%%%%%%%%%%%%%%%%%%%%%%%%%%%%%%%%%%
\begin{document}

% Here we set the space around the initial.
% Please report to http://home.gna.org/gregorio/gregoriotex/details for more details and options
\grechangedim{afterinitialshift}{2.2mm}{scalable}
\grechangedim{beforeinitialshift}{2.2mm}{scalable}
\grechangedim{interwordspacetext}{0.22 cm plus 0.15 cm minus 0.05 cm}{scalable}%
\grechangedim{annotationraise}{-0.2cm}{scalable}

% Here we set the initial font. Change 38 if you want a bigger initial.
% Emit the initials in red.
\grechangestyle{initial}{\color{red}\fontsize{38}{38}\selectfont}

\pagestyle{empty}

%%%% Titulni stranka
\begin{titulusOfficii}
\ifx\titulus\undefined
\nomenFesti{Feria IV \hebdomada{}}
\else
\titulus
\fi
\end{titulusOfficii}

\vfill

\begin{center}
%Ad usum et secundum consuetudines chori \guillemotright{}Conventus Choralis\guillemotleft.

%Editio Sancti Wolfgangi \annusEditionis
\end{center}

\scriptura{}

\pars{}

\pagebreak

\renewcommand{\headrulewidth}{0pt} % no horiz. rule at the header
\fancyhf{}
\pagestyle{fancy}

\cantusSineNeumas

\ifx\oratio\undefined
\ifx\lauda\undefined
\else
\newcommand{\oratio}{\pars{Oratio.}

\noindent Exáudi nos, Deus, salutáris noster et nos pérfice sectatóres lucis et operários veritátis, ut, qui ex te nati sumus lucis fílii, tui testes coram homínibus esse valeámus.

\pars{Pro pace in universo mundo.} \scriptura{Sir. 50, 25; 2 Esdr. 4, 20; \textbf{H416}}

\vspace{-4mm}

\antiphona{II D}{temporalia/ant-dapacemdomine.gtex}

\vfill

\noindent Deus, a quo sancta desidéria, recta consília et iusta sunt ópera: da servis tuis illam, quam mundus dare non potest, pacem; ut et corda nostra mandátis tuis dédita, et hóstium subláta formídine, témpora sint tua protectióne tranquílla.

\noindent Per Dóminum nostrum Iesum Christum, Fílium tuum, qui tecum vivit et regnat in unitáte Spíritus Sancti, Deus, per ómnia sǽcula sæculórum.

\noindent \Rbardot{} Amen.}
\fi
\ifx\laudb\undefined
\else
\newcommand{\oratio}{\pars{Oratio.}

\noindent Emítte, quǽsumus, Dómine, in corda nostra tui lúminis claritátem, ut, in via mandatórum tuórum iúgiter ambulántes, nihil umquam patiámur erróris.

\pars{Pro pace in universo mundo.} \scriptura{Sir. 50, 25; 2 Esdr. 4, 20; \textbf{H416}}

\vspace{-4mm}

\antiphona{II D}{temporalia/ant-dapacemdomine.gtex}

\vfill

\noindent Deus, a quo sancta desidéria, recta consília et iusta sunt ópera: da servis tuis illam, quam mundus dare non potest, pacem; ut et corda nostra mandátis tuis dédita, et hóstium subláta formídine, témpora sint tua protectióne tranquílla.

\noindent Per Dóminum nostrum Iesum Christum, Fílium tuum, qui tecum vivit et regnat in unitáte Spíritus Sancti, Deus, per ómnia sǽcula sæculórum.

\noindent \Rbardot{} Amen.}
\fi
\ifx\laudc\undefined
\else
\newcommand{\oratio}{\pars{Oratio.}

\noindent Sénsibus nostris, quǽsumus, Dómine, lumen sanctum tuum benígnus infúnde, ut tibi semper simus conversatióne devóti, cuius sapiéntia creáti sumus et providéntia gubernámur.

\pars{Pro pace in universo mundo.} \scriptura{Sir. 50, 25; 2 Esdr. 4, 20; \textbf{H416}}

\vspace{-4mm}

\antiphona{II D}{temporalia/ant-dapacemdomine.gtex}

\vfill

\noindent Deus, a quo sancta desidéria, recta consília et iusta sunt ópera: da servis tuis illam, quam mundus dare non potest, pacem; ut et corda nostra mandátis tuis dédita, et hóstium subláta formídine, témpora sint tua protectióne tranquílla.

\noindent Per Dóminum nostrum Iesum Christum, Fílium tuum, qui tecum vivit et regnat in unitáte Spíritus Sancti, Deus, per ómnia sǽcula sæculórum.

\noindent \Rbardot{} Amen.}
\fi
\ifx\laudd\undefined
\else
\newcommand{\oratio}{\pars{Oratio.}

\noindent Memoráre, Dómine, testaménti tui sancti, quod sanguis Agni novo fœ́dere consecrávit, ut pópulus tuus et remissiónem obtíneat peccatórum et ad redemptiónis indesinénter profíciat increméntum.

\pars{Pro pace in universo mundo.} \scriptura{Sir. 50, 25; 2 Esdr. 4, 20; \textbf{H416}}

\vspace{-4mm}

\antiphona{II D}{temporalia/ant-dapacemdomine.gtex}

\vfill

\noindent Deus, a quo sancta desidéria, recta consília et iusta sunt ópera: da servis tuis illam, quam mundus dare non potest, pacem; ut et corda nostra mandátis tuis dédita, et hóstium subláta formídine, témpora sint tua protectióne tranquílla.

\noindent Per Dóminum nostrum Iesum Christum, Fílium tuum, qui tecum vivit et regnat in unitáte Spíritus Sancti, Deus, per ómnia sǽcula sæculórum.

\noindent \Rbardot{} Amen.}
\fi
\fi

\hora{Ad Matutinum.} %%%%%%%%%%%%%%%%%%%%%%%%%%%%%%%%%%%%%%%%%%%%%%%%%%%%%

\vspace{2mm}

\cuminitiali{}{temporalia/dominelabiamea.gtex}

\vfill
%\pagebreak

\vspace{2mm}

\ifx\invitatorium\undefined
\pars{Invitatorium.} \scriptura{Ps. 94, 6; Psalmus 94; \textbf{H451}}

\vspace{-4mm}

\antiphona{VI}{temporalia/inv-dominumquifecit.gtex}
\else
\invitatorium
\fi

\vfill
\pagebreak

\ifx\hymnusmatutinum\undefined
\ifx\hiemalis\undefined
\ifx\matua\undefined
\else
\pars{Hymnus}

\cuminitiali{II}{temporalia/hym-ScientiarumDomino-MMMA.gtex}
\fi
\ifx\matub\undefined
\else
\pars{Hymnus.}

\antiphona{VIII}{temporalia/hym-ChristeLuxVera-kempten.gtex}
\fi
\ifx\matuc\undefined
\else
\pars{Hymnus}

\cuminitiali{IV}{temporalia/hym-ScientiarumDomino-kempten.gtex}
\fi
\ifx\matud\undefined
\else
\pars{Hymnus.}

\antiphona{I}{temporalia/hym-ChristeLuxVera.gtex}
\fi
\else
\ifx\matuac\undefined
\else
\pars{Hymnus}

\cuminitiali{IV}{temporalia/hym-RerumCreator.gtex}
\fi
\ifx\matubd\undefined
\else
\pars{Hymnus.}

{
\grechangedim{interwordspacetext}{0.10 cm plus 0.15 cm minus 0.05 cm}{scalable}%
\antiphona{I}{temporalia/hym-OSatorRerum.gtex}
\grechangedim{interwordspacetext}{0.22 cm plus 0.15 cm minus 0.05 cm}{scalable}%
}
\fi
\fi
\else
\hymnusmatutinum
\fi

\vspace{-3mm}

\vfill
\pagebreak

\ifx\matutinum\undefined
\ifx\matua\undefined
\else
% MAT A
\pars{Psalmus 1.} \scriptura{Ps. 17, 2}

\vspace{-4mm}

\antiphona{VI F}{temporalia/ant-diligamtedomine.gtex}

%\vspace{-2mm}

\scriptura{Ps. 17, 2-7}

%\vspace{-2mm}

\initiumpsalmi{temporalia/ps17ii_vii-initium-vi-F-auto.gtex}

\input{temporalia/ps17ii_vii-vi-F.tex} \Abardot{}

\vfill
\pagebreak

\pars{Psalmus 2.} \scriptura{Ps. 29, 11; \textbf{H151}}

\vspace{-4mm}

\antiphona{I g}{temporalia/ant-factusestadiutormeus.gtex}

%\vspace{-2mm}

\scriptura{Ps. 17, 8-20}

%\vspace{-2mm}

\initiumpsalmi{temporalia/ps17viii_xx-initium-i-g-auto.gtex}

\input{temporalia/ps17viii_xx-i-g.tex}

\vfill

\antiphona{}{temporalia/ant-factusestadiutormeus.gtex}

\vfill
\pagebreak

\pars{Psalmus 3.} \scriptura{Ps. 17, 21}

\vspace{-4mm}

\antiphona{IV* e}{temporalia/ant-retribuetmihi.gtex}

%\vspace{-2mm}

\scriptura{Ps. 17, 21-30}

%\vspace{-2mm}

\initiumpsalmi{temporalia/ps17xxi_xxx-initium-iv_-e-auto.gtex}

\input{temporalia/ps17xxi_xxx-iv_-e.tex} \Abardot{}

\vfill
\pagebreak
\fi
\ifx\matub\undefined
\else
% MAT B
\pars{Psalmus 1.} \scriptura{Ps. 38, 2; \textbf{H93}}

\vspace{-4mm}

\antiphona{IV E}{temporalia/ant-utnondelinquam.gtex}

%\vspace{-2mm}

\scriptura{Ps. 38, 2-7}

%\vspace{-2mm}

\initiumpsalmi{temporalia/ps38i-initium-iv-E-auto.gtex}

\input{temporalia/ps38i-iv-E.tex} \Abardot{}

\vfill
\pagebreak

\pars{Psalmus 2.} \scriptura{Ier. 17, 17; \textbf{H177}}

\vspace{-4mm}

\antiphona{VII c}{temporalia/ant-nonsismihi.gtex}

%\vspace{-2mm}

\scriptura{Ps. 38, 8-14}

%\vspace{-2mm}

\initiumpsalmi{temporalia/ps38ii-initium-vii-c-auto.gtex}

\input{temporalia/ps38ii-vii-c.tex} \Abardot{}

\vfill
\pagebreak

\pars{Psalmus 3.}

\vspace{-4mm}

\antiphona{IV* e}{temporalia/ant-exspectabonomentuum.gtex}

%\vspace{-2mm}

\scriptura{Ps. 51}

%\vspace{-2mm}

\initiumpsalmi{temporalia/ps51-initium-iv_-e-auto.gtex}

\input{temporalia/ps51-iv_-e.tex} \Abardot{}

\vfill
\pagebreak
\fi
\ifx\matuc\undefined
\else
% MAT C
\pars{Psalmus 1.} \scriptura{Ps. 75, 12}

\vspace{-4mm}

\antiphona{VIII G}{temporalia/ant-deusquiglorificatur.gtex}

%\vspace{-2mm}

\scriptura{Ps. 88, 2-19}

%\vspace{-2mm}

\initiumpsalmi{temporalia/ps88ii_xix-initium-viii-G-auto.gtex}

\input{temporalia/ps88ii_xix-viii-G.tex}

\vfill

\antiphona{}{temporalia/ant-deusquiglorificatur.gtex}

\vfill
\pagebreak

\pars{Psalmus 2.} \scriptura{Ps. 131, 11; \textbf{H52}}

\vspace{-5mm}

\antiphona{VIII G}{temporalia/ant-defructuventris.gtex}

\vspace{-2mm}

\scriptura{Ps. 88, 20-30}

\vspace{-2mm}

\initiumpsalmi{temporalia/ps88xx_xxx-initium-viii-G-auto.gtex}

\input{temporalia/ps88xx_xxx-viii-G.tex} \Abardot{}

\vfill
\pagebreak

\pars{Psalmus 3.} \scriptura{Ps. 111, 2; \textbf{H365}}

\vspace{-4mm}

\antiphona{VIII G*}{temporalia/ant-potensinterra.gtex}

%\vspace{-2mm}

\scriptura{Ps. 88, 31-38}

%\vspace{-2mm}

\initiumpsalmi{temporalia/ps88xxxi_xxxviii-initium-viii-G_.gtex}

\input{temporalia/ps88xxxi_xxxviii-viii-G.tex} \Abardot{}

\vfill
\pagebreak
\fi
\ifx\matud\undefined
\else
% MAT D
\pars{Psalmus 1.} \scriptura{Ps. 102, 1; \textbf{H99}}

\vspace{-4mm}

\antiphona{VIII c}{temporalia/ant-benedicanimamea.gtex}

%\vspace{-2mm}

\scriptura{Ps. 102, 1-7}

%\vspace{-2mm}

\initiumpsalmi{temporalia/ps102i-initium-viii-C-auto.gtex}

\input{temporalia/ps102i-viii-C.tex} \Abardot{}

\vfill
\pagebreak

\pars{Psalmus 2.} \scriptura{Ps. 102, 11}

\vspace{-4mm}

\antiphona{I d}{temporalia/ant-supertimentesdominum.gtex}

%\vspace{-2mm}

\scriptura{Ps. 102, 8-16}

%\vspace{-2mm}

\initiumpsalmi{temporalia/ps102ii-initium-i-d-auto.gtex}

\input{temporalia/ps102ii-i-d.tex} \Abardot{}

\vfill
\pagebreak

\pars{Psalmus 3.} \scriptura{Ps. 102, 20; \textbf{H332}}

\vspace{-4mm}

\antiphona{III g (VIII*)}{temporalia/ant-benedicitedomino.gtex}

%\vspace{-5mm}

\scriptura{Ps. 102, 17-22}

%\vspace{-2mm}

\initiumpsalmi{temporalia/ps102iii-initium-iii-g.gtex}

\input{temporalia/ps102iii-iii-g.tex} \Abardot{}

\vfill
\pagebreak
\fi
\else
\matutinum
\fi

\pars{Versus.}

\ifx\matversus\undefined
\ifx\matua\undefined
\else
\noindent \Vbardot{} Omnes mirabántur in verbis grátiæ.

\noindent \Rbardot{} Quæ procedébant de ore ipsíus.
\fi
\ifx\matub\undefined
\else
\noindent \Vbardot{} Sustínuit ánima mea in verbo eius.

\noindent \Rbardot{} Sperávit ánima mea in Dómino.
\fi
\ifx\matuc\undefined
\else
\noindent \Vbardot{} Declarátio sermónum tuórum illúminat.

\noindent \Rbardot{} Et intelléctum dat párvulis.
\fi
\ifx\matud\undefined
\else
\noindent \Vbardot{} Viam mandatórum tuórum, Dómine, fac me intellégere.

\noindent \Rbardot{} Et exercébor in mirabílibus tuis.
\fi
\else
\matversus
\fi

\ifx\sineabsolutio\undefined
\vspace{5mm}

\sineinitiali{temporalia/oratiodominica-mat.gtex}

\vspace{5mm}

\pars{Absolutio.}

\ifx\absolutio\undefined
\cuminitiali{}{temporalia/absolutio-avinculis.gtex}
\else
\absolutio
\fi

\vfill
\pagebreak

\ifx\benedictioi\undefined
\cuminitiali{}{temporalia/benedictio-solemn-ille.gtex}
\else
\benedictioi
\fi
\fi

\vspace{7mm}

\lectioi

\ifx\sineabsolutio\undefined
\noindent \Vbardot{} Tu autem, Dómine, miserére nobis.
\noindent \Rbardot{} Deo grátias.
\fi

\vfill
\pagebreak

\responsoriumi

\vfill
\pagebreak

\ifx\sineabsolutio\undefined
\ifx\benedictioii\undefined
\cuminitiali{}{temporalia/benedictio-solemn-divinum.gtex}
\else
\benedictioii
\fi
\fi

\vspace{7mm}

\lectioii

\ifx\sineabsolutio\undefined
\noindent \Vbardot{} Tu autem, Dómine, miserére nobis.
\noindent \Rbardot{} Deo grátias.
\fi

\vfill
\pagebreak

\responsoriumii

\vfill
\pagebreak

\ifx\sineabsolutio\undefined
\ifx\benedictioiii\undefined
\cuminitiali{}{temporalia/benedictio-solemn-adsocietatem.gtex}
\else
\benedictioiii
\fi
\fi

\vspace{7mm}

\lectioiii

\ifx\sineabsolutio\undefined
\noindent \Vbardot{} Tu autem, Dómine, miserére nobis.
\noindent \Rbardot{} Deo grátias.
\fi

\vfill
\pagebreak

\responsoriumiii

\vfill
\pagebreak

\rubrica{Reliqua omittuntur, nisi Laudes separandæ sint.}

\sineinitiali{temporalia/domineexaudi.gtex}

\vfill

\oratio

\vfill

\noindent \Vbardot{} Dómine, exáudi oratiónem meam.
\Rbardot{} Et clamor meus ad te véniat.

\vfill

\noindent \Vbardot{} Benedicámus Dómino.
\noindent \Rbardot{} Deo grátias.

\vfill

\noindent \Vbardot{} Fidélium ánimæ per misericórdiam Dei requiéscant in pace.
\Rbardot{} Amen.

\vfill
\pagebreak

\hora{Ad Laudes.} %%%%%%%%%%%%%%%%%%%%%%%%%%%%%%%%%%%%%%%%%%%%%%%%%%%%%

\cantusSineNeumas

\vspace{0.5cm}
\ifx\deusinadiutorium\undefined
\grechangedim{interwordspacetext}{0.18 cm plus 0.15 cm minus 0.05 cm}{scalable}%
\cuminitiali{}{temporalia/deusinadiutorium-communis.gtex}
\grechangedim{interwordspacetext}{0.22 cm plus 0.15 cm minus 0.05 cm}{scalable}%
\else
\deusinadiutorium
\fi

\vfill
\pagebreak

\ifx\hymnuslaudes\undefined
\ifx\lauda\undefined
\else
\pars{Hymnus} \scriptura{Prudentius (\olddag{} 413)}

\cuminitiali{I}{temporalia/hym-NoxEtTenebrae.gtex}
\fi
\ifx\laudb\undefined
\else
\pars{Hymnus}

\cuminitiali{I}{temporalia/hym-FulgentisAuctor-hk.gtex}
\fi
\ifx\laudc\undefined
\else
\pars{Hymnus} \scriptura{Prudentius (\olddag{} 413)}

\cuminitiali{VIII}{temporalia/hym-NoxEtTenebrae-einsiedeln.gtex}
\fi
\ifx\laudd\undefined
\else
\pars{Hymnus}

\grechangedim{interwordspacetext}{0.16 cm plus 0.15 cm minus 0.05 cm}{scalable}%
\cuminitiali{IV}{temporalia/hym-FulgentisAuctor.gtex}
\grechangedim{interwordspacetext}{0.22 cm plus 0.15 cm minus 0.05 cm}{scalable}%
\vspace{-3mm}
\fi
\else
\hymnuslaudes
\fi

\vfill
\pagebreak

\ifx\laudes\undefined
\ifx\lauda\undefined
\else
\pars{Psalmus 1.} \scriptura{Ps. 35, 6}

\vspace{-4mm}

\antiphona{II D}{temporalia/ant-domineincaelo.gtex}

\scriptura{Psalmus 35.}

\initiumpsalmi{temporalia/ps35-initium-ii-D-auto.gtex}

\input{temporalia/ps35-ii-D.tex} \Abardot{}

\vfill
\pagebreak

\pars{Psalmus 2.} \scriptura{Idt. 16, 16}

\vspace{-4mm}

\antiphona{II* d}{temporalia/ant-dominemagnuses.gtex}

\scriptura{Canticum Iudith, Idt. 16, 2.15-19}

\initiumpsalmi{temporalia/iudith2-initium-ii_-d-auto.gtex}

\input{temporalia/iudith2-ii_-d.tex} \Abardot{}

\vfill
\pagebreak

\pars{Psalmus 3.} \scriptura{Ps. 46, 7-8; \textbf{H44}}

\vspace{-4mm}

\antiphona{VIII c}{temporalia/ant-rexomnisterrae.gtex}

%\vspace{-2mm}

\scriptura{Psalmus 46.}

%\vspace{-2mm}

\initiumpsalmi{temporalia/ps46-initium-viii-c-auto.gtex}

\input{temporalia/ps46-viii-c.tex} \Abardot{}

\vfill
\pagebreak
\fi
\ifx\laudb\undefined
\else
\pars{Psalmus 1.} \scriptura{Ps. 76, 14}

\vspace{-4mm}

\antiphona{VIII G}{temporalia/ant-deusinsancto.gtex}

\scriptura{Psalmus 76.}

\initiumpsalmi{temporalia/ps76-initium-viii-G-auto.gtex}

\input{temporalia/ps76-viii-G.tex}

\vfill

\antiphona{}{temporalia/ant-deusinsancto.gtex}

\vfill
\pagebreak

\pars{Psalmus 2.} \scriptura{1 Reg. 2, 1; \textbf{H96}}

\vspace{-4mm}

\antiphona{II D}{temporalia/ant-exsultavitcormeum.gtex}

\vspace{-4mm}

\scriptura{Canticum Annæ, 1 Reg. 2, 1-10}

\vspace{-3mm}

\initiumpsalmi{temporalia/anna-initium-ii-D-auto.gtex}

\input{temporalia/anna-ii-D.tex}

\vfill

\antiphona{}{temporalia/ant-exsultavitcormeum.gtex}

\vfill
\pagebreak

\pars{Psalmus 3.} \scriptura{Ps. 96, 1}

\vspace{-4mm}

\antiphona{III a\textsuperscript{2}}{temporalia/ant-dominusregnavit.gtex}

\vspace{-2mm}

\scriptura{Psalmus 96.}

\vspace{-2mm}

\initiumpsalmi{temporalia/ps96-initium-iii-a2.gtex}

\input{temporalia/ps96-iii-a2.tex} \Abardot{}

\vfill
\pagebreak
\fi
\ifx\laudc\undefined
\else
\pars{Psalmus 1.} \scriptura{Ps. 85, 1}

\vspace{-4mm}

\antiphona{II D}{temporalia/ant-inclinadomineaurem.gtex}

\scriptura{Psalmus 85.}

\initiumpsalmi{temporalia/ps85-initium-ii-D-auto.gtex}

\input{temporalia/ps85-ii-D.tex}

\vfill

\antiphona{}{temporalia/ant-inclinadomineaurem.gtex}

\vfill
\pagebreak

\pars{Psalmus 2.} \scriptura{Ps. 118, 1; \textbf{H91}}

\vspace{-4mm}

\antiphona{VIII G}{temporalia/ant-beatiquiambulant.gtex}

%\vspace{-2mm}

\scriptura{Canticum Isaiaæ, Is. 33, 13-16}

%\vspace{-2mm}

\initiumpsalmi{temporalia/isaiae8-initium-viii-g-auto.gtex}

\input{temporalia/isaiae8-viii-g.tex} \Abardot{}

\vfill
\pagebreak

\pars{Psalmus 3.} \scriptura{Ps. 97, 1; \textbf{H98}}

\vspace{-4mm}

\antiphona{E}{temporalia/ant-quiamirabiliafecit.gtex}

%\vspace{-2mm}

\scriptura{Psalmus 97.}

%\vspace{-2mm}

\initiumpsalmi{temporalia/ps97-initium-e.gtex}

\input{temporalia/ps97-e.tex} \Abardot{}

\vfill
\pagebreak
\fi
\ifx\laudd\undefined
\else
\pars{Psalmus 1.} \scriptura{Ps. 107, 2}

\vspace{-4mm}

\antiphona{VIII G}{temporalia/ant-paratumcormeum.gtex}

\vspace{-2mm}

\scriptura{Psalmus 107.}

\vspace{-2mm}

\initiumpsalmi{temporalia/ps107-initium-viii-G-auto.gtex}

\input{temporalia/ps107-viii-G.tex} \Abardot{}

\vfill
\pagebreak

\pars{Psalmus 2.} \scriptura{Is. 61, 10}

\vspace{-4mm}

\antiphona{VII c}{temporalia/ant-induitmedominus.gtex}

\vspace{-1mm}

\scriptura{Canticum Isaiaæ, Is. 61, 10-11; 62, 1-7}

\vspace{-3mm}

\initiumpsalmi{temporalia/isaiae4-initium-vii-c-auto.gtex}

\vspace{-1.5mm}

\input{temporalia/isaiae4-vii-c.tex} \Abardot{}

\vfill
\pagebreak

\pars{Psalmus 3.} \scriptura{Ps. 145, 2; \textbf{H100}}

\vspace{-4mm}

\antiphona{IV* e}{temporalia/ant-laudabodeum.gtex}

\scriptura{Psalmus 145.}

\initiumpsalmi{temporalia/ps145-initium-iv_-e-auto.gtex}

\input{temporalia/ps145-iv_-e.tex} \Abardot{}

\vfill
\pagebreak
\fi
\else
\laudes
\fi

\ifx\lectiobrevis\undefined
\ifx\lauda\undefined
\else
\pars{Lectio Brevis.} \scriptura{Tob. 4, 14-15.16.19}

\noindent Atténde tibi, fili, in ómnibus opéribus tuis et esto sápiens in ómnibus sermónibus tuis et, quod óderis, némini féceris. De pane tuo commúnica esuriénti et de vestiméntis tuis nudis; ex ómnibus, quæcúmque tibi abundáverint, fac eleemósynam. Omni témpore bénedic Dóminum et póstula ab illo, ut dirigántur viæ tuæ et omnes sémitæ tuæ et consília bene disponántur.
\fi
\ifx\laudb\undefined
\else
\pars{Lectio Brevis.} \scriptura{Rom. 8, 35.37}

\noindent Quis nos separábit a caritáte Christi? Tribulátio an angústia an persecútio an fames an núditas an perículum an gládius? Sed in his ómnibus supervíncimus per eum, qui diléxit nos.
\fi
\ifx\laudc\undefined
\else
\pars{Lectio Brevis.} \scriptura{Iob 1, 21; 2, 10}

\noindent Nudus egréssus sum de útero matris meæ et nudus revértar illuc. Dóminus dedit, Dóminus ábstulit; sicut Dómino plácuit, ita factum est: sit nomen Dómini benedíctum. Si bona suscépimus de manu Dei, mala quare non suscipiámus?
\fi
\ifx\laudd\undefined
\else
\pars{Lectio Brevis.} \scriptura{Dt. 4, 39-40}

\noindent Scito hódie et cogitáto in corde tuo quod Dóminus ipse sit Deus in cælo sursum et in terra deórsum, et non sit álius. Custódi præcépta eius atque mandáta, quæ ego præcípio tibi hódie.
\fi
\else
\lectiobrevis
\fi

\vfill

\ifx\responsoriumbreve\undefined
\ifx\laudac\undefined
\else
\pars{Responsorium breve.} \scriptura{Ps. 118, 36-37}

\cuminitiali{E}{temporalia/resp-inclinacormeumdeus.gtex}
\fi
\ifx\laudbd\undefined
\else
\pars{Responsorium breve.} \scriptura{Ps. 33, 2}

\cuminitiali{VI}{temporalia/resp-benedicamdominum.gtex}
\fi
\else
\responsoriumbreve
\fi
\vfill
\pagebreak

\ifx\benedictus\undefined
\ifx\laudac\undefined
\else
\pars{Canticum Zachariæ.} \scriptura{Lc. 1, 72.68}

\vspace{-4mm}

{
\grechangedim{interwordspacetext}{0.18 cm plus 0.15 cm minus 0.05 cm}{scalable}%
\antiphona{VIII G}{temporalia/ant-memoraredominetestamenti.gtex}
\grechangedim{interwordspacetext}{0.22 cm plus 0.15 cm minus 0.05 cm}{scalable}%
}

%\vspace{-3mm}

\scriptura{Lc. 1, 68-79}

%\vspace{-1mm}

\initiumpsalmi{temporalia/benedictus-initium-viii-G-auto.gtex}

\input{temporalia/benedictus-viii-G.tex} \Abardot{}
\fi
\ifx\laudbd\undefined
\else
\pars{Canticum Zachariæ.} \scriptura{Lc. 1, 74-75}

\vspace{-4mm}

{
\grechangedim{interwordspacetext}{0.18 cm plus 0.15 cm minus 0.05 cm}{scalable}%
\antiphona{VIII G\textsuperscript{2}}{temporalia/ant-liberatiserviamusdeo.gtex}
\grechangedim{interwordspacetext}{0.22 cm plus 0.15 cm minus 0.05 cm}{scalable}%
}

%\vspace{-3mm}

\scriptura{Lc. 1, 68-79}

%\vspace{-1mm}

\initiumpsalmi{temporalia/benedictus-initium-viii-G2-auto.gtex}

\input{temporalia/benedictus-viii-G2.tex} \Abardot{}
\fi
\else
\benedictus
\fi

\vspace{-1cm}

\vfill
\pagebreak

\ifx\precestotum\undefined
\pars{Preces.}

\sineinitiali{}{temporalia/tonusprecumnovum.gtex}

\ifx\preces\undefined
\ifx\lauda\undefined
\else
\noindent Grátias agámus Christo eúmque semper laudémus,~\gredagger{} quia non dedignátur fratres vocáre quos sanctíficat.~\grestar{} Ideo ei supplicémus:

\Rbardot{} Sanctífica fratres tuos, Dómine.

\noindent Fac ut huius diéi inítia in honórem resurrectiónis tuæ puris tibi córdibus consecrémus,~\grestar{} et diem totum tibi gratum sanctificatiónis opéribus faciámus.

\Rbardot{} Sanctífica fratres tuos, Dómine.

\noindent Qui diem, amóris tui signum, ad salútem et lætítiam nobis renovásti,~\grestar{} rénova nos cotídie ad glóriam tuam.

\Rbardot{} Sanctífica fratres tuos, Dómine.

\noindent Doce nos hódie te in ómnibus præséntem agnóscere,~\grestar{} teque in mæréntibus præsértim et paupéribus inveníre.

\Rbardot{} Sanctífica fratres tuos, Dómine.

\noindent Da nos hódie cum ómnibus pacem habére,~\grestar{} némini vero malum réddere pro malo.

\Rbardot{} Sanctífica fratres tuos, Dómine.
\fi
\ifx\laudb\undefined
\else
\noindent Benedíctus Deus salvátor noster,~\gredagger{} qui usque ad consummatiónem sǽculi se ómnibus diébus cum Ecclésia sua mansúrum promísit.~\grestar{} Ideo ei grátias agéntes clamémus:

\Rbardot{} Mane nobíscum, Dómine.

\noindent Mane nobíscum, Dómine, toto die,~\grestar{} numquam declínet a nobis sol grátiæ tuæ.

\Rbardot{} Mane nobíscum, Dómine.

\noindent Hunc diem tibi tamquam oblatiónem consecrámus,~\grestar{} dum nos nihil pravi factúros aut probatúros pollicémur.

\Rbardot{} Mane nobíscum, Dómine.

\noindent Fac, Dómine, ut donum lucis hic totus dies evádat,~\grestar{} ut simus sal terræ et lux mundi.

\Rbardot{} Mane nobíscum, Dómine.

\noindent Spíritus Sancti tui cáritas dírigat corda et lábia nostra,~\grestar{} ut in tua iustítia semper et laude maneámus.

\Rbardot{} Mane nobíscum, Dómine.
\fi
\ifx\laudc\undefined
\else
\noindent Christum, qui nutrit et fovet Ecclésiam pro qua seípsum trádidit,~\grestar{} hac deprecatióne rogémus:

\Rbardot{} Réspice Ecclésiam tuam, Dómine.

\noindent Pastor Ecclésiæ tuæ, benedíctus es,~\gredagger{} qui nobis hódie tríbuis lucem et vitam,~\grestar{} nos gratos redde pro múnere tam iucúndo.

\Rbardot{} Réspice Ecclésiam tuam, Dómine.

\noindent Gregem tuum misericórditer réspice,~\gredagger{} quem in tuo nómine congregásti,~\grestar{} ne quis péreat eórum, quos Pater dedit tibi.

\Rbardot{} Réspice Ecclésiam tuam, Dómine.

\noindent Ecclésiam tuam duc in via mandatórum tuórum,~\grestar{} fidélem tibi Spíritus Sanctus eam effíciat.

\Rbardot{} Réspice Ecclésiam tuam, Dómine.

\noindent Ecclésiam ad mensam verbi tui panísque vivífica,~\grestar{} ut in fortitúdine huius cibi te læta sequátur.

\Rbardot{} Réspice Ecclésiam tuam, Dómine.
\fi
\ifx\laudd\undefined
\else
\noindent Christus, splendor patérnæ glóriæ,~\gredagger{} verbo suo nos illúminat.~\grestar{} Eum amánter invocémus dicéntes:

\Rbardot{} Exáudi nos, Rex ætérnæ glóriæ.

\noindent Benedíctus es, auctor fídei nostræ et consummátor,~\grestar{} qui de ténebris vocásti nos in admirábile lumen tuum.

\Rbardot{} Exáudi nos, Rex ætérnæ glóriæ.

\noindent Qui cæcórum óculos aperuísti,~\gredagger{} et surdos fecísti audíre,~\grestar{} ádiuva incredulitátem nostram.

\Rbardot{} Exáudi nos, Rex ætérnæ glóriæ.

\noindent Dómine, in dilectióne tua iúgiter maneámus,~\grestar{} ne ab ínvicem separémur.

\Rbardot{} Exáudi nos, Rex ætérnæ glóriæ.

\noindent Da nobis in tentatióne resístere,~\gredagger{} in tribulatióne sustinére,~\grestar{} et in prósperis grátias ágere.

\Rbardot{} Exáudi nos, Rex ætérnæ glóriæ.
\fi
\else
\preces
\fi

\vfill

\pars{Oratio Dominica.}

\cuminitiali{}{temporalia/oratiodominicaalt.gtex}

\vfill
\pagebreak

\pars{Deprecatio Gelasii}

\vspace{-5mm}

\grechangedim{interwordspacetext}{0.16 cm plus 0.15 cm minus 0.05 cm}{scalable}%
\antiphona{D\textsuperscript{1}}{temporalia/deprecatio4-propace.gtex}
\grechangedim{interwordspacetext}{0.22 cm plus 0.15 cm minus 0.05 cm}{scalable}%

\vfill

\pars{Oratio Dominica.}

\cuminitiali{D}{temporalia/oratiodominica-d.gtex}
\else
\precestotum
\fi

\vfill
\pagebreak

% Oratio. %%%
\oratio

\vspace{-1mm}

\vfill

\rubrica{Hebdomadarius dicit Dominus vobiscum, vel, absente sacerdote vel diacono, sic concluditur:}

\vspace{2mm}

\ifx\dominusnosbenedicat\undefined
\antiphona{C}{temporalia/dominusnosbenedicat.gtex}
\else
\dominusnosbenedicat
\fi

\rubrica{Postea cantatur a cantore:}

\vspace{2mm}

\ifx\benedicamuslaudes\undefined
\cuminitiali{IV}{temporalia/benedicamus-feria-laudes.gtex}
\else
\benedicamuslaudes
\fi

\vspace{1mm}

\vfill
\pagebreak

\end{document}

